\section{The interaction shape and its three-state metric sector}
\label{sec:shape}
The mass dilation of Section~\ref{sec:quartic} leaves the chiral mass
class zero. A genuine additional parameter is the dimensionless shape
$c\in\C^*$ in
\begin{equation}\label{shape:potential}
G_c=\frac{A^2+B^2}{2}+\frac c4(A^4+B^4+A^2B^2),
\qquad f_c=G_c/4,\qquad D_c=\bar\partial+df_c\wedge .
\end{equation}
The physical family is defined on all of $\C^*$; the real-plane period
below has the smaller directly convergent domain $\re c>0$.

\subsection{Jacobi multiplication, symmetry and residue}
Set $s=A^2+B^2$, $t=A^2B^2$, and $\Omega=dA\wedge dB$. The Jacobi
ideal is generated by
$A(1+cA^2+cB^2/2)$ and $B(1+cB^2+cA^2/2)$. Its nine simple points
are the origin, the four axis points with nonzero coordinate squared
$-1/c$, and the four points with $A^2=B^2=-2/(3c)$. A basis is
\[
1,A,B,A^2,AB,B^2,A^2B,AB^2,A^2B^2.
\]
The invariant subalgebra has the exact relations
\begin{equation}\label{shape:algebra}
s^2=-s/c+t,\qquad st=-4t/(3c),\qquad t^2=4t/(9c^2).
\end{equation}
They follow by reducing the cubic gradient relations; evaluation at
the nine points proves independence of the displayed basis.

The group $(\mathbb Z_2)^2\rtimes S_2$ acts by sign changes and exchange.
The volume $\Omega$ transforms by its determinant character. The
corresponding physical symmetry sector therefore consists of invariant
polynomials times $\Omega$, and has basis $(1,s,t)\Omega$. The
Hamiltonian, harmonic projection and shape insertion preserve this
sector. It is a three-dimensional block of the nine-vacuum theory;
there is no division of physical norms by the group order eight.

The canonical Higgs insertion is multiplication by
\begin{equation}\label{shape:higgs}
[\partial_cf_c]=[ (A^4+B^4+A^2B^2)/16]=-s/(16c).
\end{equation}
It has eigenvalues $0,1/(16c^2),1/(12c^2)$, and the identity is
cyclic under it, since its first two powers generate $s$ and $s^2$.
Thus a scalar closure of this physical identity sector is insufficient.

Use the holomorphic source frame
\begin{equation}\label{shape:frame}
E=(1,r,v)\Omega,\qquad r=cs,\quad v=c^2t.
\end{equation}
In this frame the multiplication matrix of $r$ and the residue matrix are
\begin{equation}\label{shape:matrices}
R=\begin{pmatrix}0&0&0\\1&-1&0\\0&1&-4/3\end{pmatrix},
\qquad
\eta_E=\frac{64}{3}
\begin{pmatrix}0&0&1\\0&1&-4/3\\1&-4/3&4/9\end{pmatrix}.
\end{equation}
Here the residue convention is
\[
\eta(p\Omega,q\Omega)=\frac1{(2\pi i)^2}\int
\frac{pq\,dA\wedge dB}{(f_c)_A(f_c)_B}
=\sum_{df_c(a)=0}\frac{p(a)q(a)}{\det\operatorname{Hess}f_c(a)} .
\]
The inverse Hessian determinants per point are $16,-16,12$ on the
three types, so the orbit weights are $(16,-64,48)$. This proves
\eqref{shape:matrices}, including $\eta(1,1)=0$ and
$\eta_E R=R^T\eta_E$. The orbit idempotents are
\begin{equation}\label{shape:idempotents}
e_0=1+r+\tfrac34v,\qquad e_a=-r-3v,\qquad e_b=\tfrac94v.
\end{equation}
Their residue is $\operatorname{diag}(16,-64,48)$. A zero identity
self-residue does not imply a zero Hermitian norm.

\subsection{Radiality and the conditional chiral equation}
Let $g_E(c,\bar c)$ be the physical metric of the fixed sources
\eqref{shape:frame}. Dilation gives
\[
f_c=\phi_c^*((1/c)f_1),\qquad
\phi_c(A,B)=(\sqrt c\,A,\sqrt c\,B).
\]
All three sources become $c^{-1}$ times pullbacks of fixed polynomial
volume sources. Middle-degree pullback preserves the $L^2$ norm
under a constant conformal dilation in real dimension four. For the
reference differential $\bar\partial+u\,df_1\wedge$, multiplication
of a $(p,q)$-form by $e^{ip\theta}$ conjugates $u$ to $e^{i\theta}u$.
All three volume sources have the same holomorphic degree. Their
fixed-source Gram matrix is therefore invariant under this phase.
Consequently the entire physical matrix in \eqref{shape:frame} is
radial, not only its identity component.

Put $u=1/c$, $\rho=|u|$, and write this metric as $g(\rho)$. Under
the \emph{additional identification} of the physical Hodge family
with the exact chiral equation in this source frame, its equations are
\begin{align}
\partial_{\bar u}(g^{-1}\partial_u g)
&=[C,g^{-1}C^\dagger g],\qquad C=R/16,\label{shape:ttstar}\\
(g^{-1}g')'+\rho^{-1}g^{-1}g'
&=\tfrac1{64}[R,g^{-1}R^Tg].\label{shape:radial}
\end{align}
The factor $1/16$ follows from $f_c=G_c/4$; the displayed physical
superpotential variation would be twice this insertion.
The Hodge and chiral frameworks are those of
\cite{Fan,CecottiVafa,SonnerTong}. We retain the exact identification
in \eqref{shape:ttstar} as an explicit hypothesis for the metric
selection theorem below; Hodge existence alone is not its proof.

The eigenvalues of $C$ are $0,-1/16,-1/12$, the three critical-value
scales of $u f_1$. Neither these scales nor the algebra fix Stokes
data. For example, in the orbit frame a constant positive metric
proportional to $\operatorname{diag}(16,64,48)$ makes $C$ normal
and has zero curvature. With a consistent overall residue normalization
it satisfies the reality relation as well. The physical endpoint in
Section~\ref{sec:endpoints} excludes this constant control in the
specified frame.

\subsection{The exact shape-period equation}
The separate holomorphic period is
\begin{equation}\label{shape:period}
\mathcal C(c)=\frac1{2\pi}\int_{\R^2}e^{-G_c(A,B)}\,dA\,dB,
\qquad \re c>0.
\end{equation}
Quartic decay justifies holomorphic differentiation on compact subsets.
For real $c>0$,
$(-1)^k\mathcal C^{(k)}(c)>0$ for every $k\ge0$, because it is the
Gaussian expectation of $( (A^4+B^4+A^2B^2)/4)^k e^{-cq_4/4}$.
In particular it is strictly decreasing and $\mathcal C(0+)=1$.

\begin{proposition}[Shape-period transport]\label{shape:period-ode}
The period \eqref{shape:period} satisfies
\begin{equation}\label{shape:ode}
\begin{split}
48c^4\mathcal C'''+(288c^3+28c^2)\mathcal C''
+(324c^2+56c+4)\mathcal C'\\
+(36c+7)\mathcal C=0.
\end{split}
\end{equation}
Its first one-sided asymptotic derivatives are
\[
\mathcal C(0+)=1,\quad \mathcal C'(0+)=-7/4,\quad
\mathcal C''(0+)=297/16,\quad \mathcal C'''(0+)=-31815/64.
\]
\end{proposition}
\begin{proof}
For the mass/coupling integral $I(M,g)$ of Section~\ref{sec:quartic},
radial integration by parts gives
$(M\partial_M+2g\partial_g)I=-I$.
Writing $A_n(c)=\partial_M^n I(1,c)$ yields
$A_{n+1}=-(n+1)A_n-2cA_n'$, hence
\[
\begin{aligned}
A_1&=-\mathcal C-2c\mathcal C',\\
A_2&=2\mathcal C+10c\mathcal C'+4c^2\mathcal C'',\\
A_3&=-6\mathcal C-54c\mathcal C'-48c^2\mathcal C''-8c^3\mathcal C'''.
\end{aligned}
\]
Insert these expressions into the fixed-quartic mass equation.
The derivative values follow independently from Gaussian moments.
Taylor's finite-order remainder for the exponential is integrably
dominated on $\re c\ge0$, so these are controlled asymptotic jets.
No convergent Taylor series through the singular point $c=0$ is asserted.
\end{proof}

\subsection{The physical-to-period ratio and scalar compensation}
Let $H(c,\bar c)=g_{00}(c,\bar c)>0$ be the raw identity-volume norm,
and define $\mathcal R=H/|\mathcal C|^2$ where the period is nonzero.
A common source rescaling cancels from this ratio.

\begin{proposition}[No constant ratio on a complex-open shape region]
\label{shape:ratio}
The ratio $\mathcal R$ is constant on no nonempty complex-open
subset of $\re c>0$. A common scalar source adjustment cannot
preserve both a fixed identity Euler period and its physical Euler
metric along a nonconstant holomorphic shape variation.
\end{proposition}
\begin{proof}
If $H=K|\mathcal C|^2$ on such a patch, radiality of $H$ gives
$\im(c\mathcal C'/\mathcal C)=0$. This holomorphic expression is
then a real constant $k$. The identity theorem extends
$c\mathcal C'-k\mathcal C=0$ throughout the right half-plane.
Along the positive axis $\mathcal C(c)=Ac^k$; its limit one at
zero forces $k=0,A=1$, contrary to strict decrease.

For a holomorphic shape $c=c(M)$, simultaneous preservation would
require $a(M)\mathcal C(c(M))=C_*$ and
$|a(M)|^2H(c(M),\overline{c(M)})=H_*$ with $C_*\ne0,H_*>0$.
Cancellation makes $\mathcal R$ constant on the image of $c(M)$,
which is open unless $c(M)$ is constant. This proves the claim;
the compensating scalar need not itself be holomorphic.
\end{proof}
Even preservation of the identity period alone with a fixed scalar
source forces $c(M)$ constant, by the open mapping theorem and
nonconstancy of $\mathcal C$. A variable scalar
$a(c)=\mathcal C(c_0)/\mathcal C(c)$ can restore that period, but
the resulting physical metric retains the factor $\mathcal R$.

The angular discrepancy can be exhibited without evaluating $H$:
\begin{equation}\label{shape:angular-ratio}
\frac{\mathcal R(re^{i\theta},re^{-i\theta})}{\mathcal R(r,r)}
=1-\tfrac72r(1-\cos\theta)+O(r^2),
\qquad |\theta|<\pi/2.
\end{equation}
It follows from $\mathcal C(c)=1-7c/4+O(|c|^2)$ and exact radiality.
This does not establish variation or monotonicity of the ratio along
the positive real axis alone. State mixing, different boundaries,
cycles or target metrics change the hypotheses and remain possible.
